\documentclass[11pt,a4paper]{article}

% ── Codificación y fuentes ────────────────────────────────────────────────────
\usepackage[utf8]{inputenc}
\usepackage[T1]{fontenc}
\usepackage{lmodern}
\usepackage[spanish]{babel}

% ── Geometría y espaciado ─────────────────────────────────────────────────────
\usepackage[top=2.5cm, bottom=2.5cm, left=2.8cm, right=2.8cm]{geometry}
\usepackage{setspace}
\setstretch{1.15}
\usepackage{parskip}

% ── Tablas ────────────────────────────────────────────────────────────────────
\usepackage{booktabs}
\usepackage{tabularx}
\usepackage{array}
\usepackage{longtable}
\usepackage{enumitem}

% ── Colores y cajas ───────────────────────────────────────────────────────────
\usepackage[dvipsnames]{xcolor}
\usepackage{tcolorbox}
\tcbuselibrary{skins, breakable}

% ── Cabeceras y pies de página ────────────────────────────────────────────────
\usepackage{fancyhdr}
\usepackage{lastpage}
\pagestyle{fancy}
\fancyhf{}
\fancyhead[L]{\small\textbf{Informe de Evaluación} --- Física para Bioanalistas}
\fancyhead[R]{\small Debora Echezuria}
\fancyfoot[C]{\small Página \thepage\ de \pageref{LastPage}}
\renewcommand{\headrulewidth}{0.4pt}

% ── Hiperenlaces ──────────────────────────────────────────────────────────────
\usepackage[colorlinks=true, linkcolor=NavyBlue, urlcolor=NavyBlue]{hyperref}

% ── Paleta de colores personalizados ─────────────────────────────────────────
\definecolor{desempeno_excelente}{HTML}{1A6B2F}
\definecolor{desempeno_bueno}{HTML}{2E5A9C}
\definecolor{desempeno_suficiente}{HTML}{8B6914}
\definecolor{desempeno_deficiente}{HTML}{C0392B}
\definecolor{desempeno_insuficiente}{HTML}{7B241C}
\definecolor{grisFondo}{HTML}{F5F5F5}
\definecolor{azulTitulo}{HTML}{1B2A6B}
\definecolor{verdeFortaleza}{HTML}{1A6B2F}
\definecolor{naranjaMejora}{HTML}{A04000}

% ── Estilos de cajas ─────────────────────────────────────────────────────────
\tcbset{
  cajaTitulo/.style={
    enhanced, breakable,
    colback=azulTitulo!8, colframe=azulTitulo,
    fonttitle=\bfseries\large, coltitle=white,
    attach boxed title to top left={yshift=-2mm, xshift=4mm},
    boxed title style={colback=azulTitulo},
    top=6pt, bottom=6pt, left=8pt, right=8pt,
  },
  cajaFortaleza/.style={
    enhanced, breakable,
    colback=verdeFortaleza!6, colframe=verdeFortaleza!60,
    fonttitle=\bfseries, coltitle=verdeFortaleza,
    top=4pt, bottom=4pt, left=8pt, right=8pt,
  },
  cajaMejora/.style={
    enhanced, breakable,
    colback=naranjaMejora!6, colframe=naranjaMejora!60,
    fonttitle=\bfseries, coltitle=naranjaMejora,
    top=4pt, bottom=4pt, left=8pt, right=8pt,
  },
  cajaRecomendacion/.style={
    enhanced, breakable,
    colback=grisFondo, colframe=gray!50,
    fonttitle=\bfseries, coltitle=black,
    top=4pt, bottom=4pt, left=8pt, right=8pt,
  },
}

% ─────────────────────────────────────────────────────────────────────────────
\begin{document}

% ══ PORTADA ══════════════════════════════════════════════════════════════════
\begin{center}
  {\Large\bfseries\color{azulTitulo} Universidad de Oriente/Núcleo Sucre --- Física para Bioanalistas}\\[4pt]
  {\large Informe de Evaluación Preliminar}\\[12pt]
  \rule{\linewidth}{1.2pt}\\[6pt]
  {\LARGE\bfseries Debora Echezuria}\\[4pt]
  \rule{\linewidth}{0.6pt}
\end{center}

% ══ FICHA TÉCNICA ════════════════════════════════════════════════════════════
\vspace{4pt}
\begin{tcolorbox}[cajaTitulo, title=Datos del informe]
\begin{tabular}{@{}llll@{}}
  \textbf{Estudiante:}  & Debora Echezuria  &
  \textbf{Fecha:}       & 2026-08-08 \\[4pt]
  \textbf{Archivo PDF:} & \texttt{Debora\_Echezuria.pdf}  &
  \textbf{Páginas:}     & 4 \\
\end{tabular}
\end{tcolorbox}

% ══ RESULTADO GLOBAL ═════════════════════════════════════════════════════════
\vspace{6pt}
\begin{center}
  \begin{tcolorbox}[
    enhanced,
    colback=white, colframe=desempeno_suficiente,
    width=0.62\linewidth,
    boxrule=2pt,
    halign=center, valign=center,
    top=8pt, bottom=8pt,
  ]
    \centering
    {\large\bfseries Puntaje sugerido}\\[6pt]
    {\Huge\bfseries\color{desempeno_suficiente} 6.4/10}\\[6pt]
    {\large Nivel de desempeño: \textbf{\color{desempeno_suficiente} Suficiente}}
  \end{tcolorbox}
\end{center}

% ══ RESUMEN GENERAL ══════════════════════════════════════════════════════════
\section*{Resumen general}
El estudiante demuestra un entendimiento básico de varios principios físicos aplicados a bioanálisis, destacando en la resolución de problemas numéricos y la identificación de fenómenos. Sin embargo, presenta deficiencias significativas en la argumentación cualitativa, la explicación detallada de los mecanismos físicos y sus implicaciones clínicas. La omisión completa de una pregunta clave sobre capilaridad impacta negativamente el desempeño general.

% ══ FORTALEZAS Y ÁREAS DE MEJORA ════════════════════════════════════════════
\vspace{4pt}
\begin{minipage}[t]{0.48\linewidth}
  \begin{tcolorbox}[cajaFortaleza, title={Fortalezas}]
\begin{itemize}[leftmargin=1.5em, itemsep=2pt]
  \item Comprensión clara del rol del calor latente de vaporización en la termorregulación por sudoración.
  \item Identificación precisa de los riesgos clínicos asociados a la alta humedad y la falta de evaporación.
  \item Aplicación correcta de la Primera Ley de Fick para calcular cambios en la eficiencia de difusión.
  \item Identificación acertada de fenómenos como hemólisis y crenación.
  \item Cálculo preciso de la presión osmótica utilizando la ecuación de van 't Hoff.
  \item Aplicación correcta de la Ley de Laplace para el cálculo de presión en alveolos.
  \item Buen manejo de unidades en la mayoría de los cálculos.
\end{itemize}
  \end{tcolorbox}
\end{minipage}
\hfill
\begin{minipage}[t]{0.48\linewidth}
  \begin{tcolorbox}[cajaMejora, title={Areas de mejora}]
\begin{itemize}[leftmargin=1.5em, itemsep=2pt]
  \item Profundizar en la explicación cualitativa de los mecanismos físicos y sus consecuencias biológicas/clínicas.
  \item Completar todas las partes de las preguntas, evitando dejar ítems en blanco.
  \item Mejorar la precisión en la descripción de procesos físicos (ej. mecanismo de ósmosis en crenación).
  \item Revisar los principios de capilaridad y su aplicación en contextos biológicos.
\end{itemize}
  \end{tcolorbox}
\end{minipage}

% ══ OBSERVACIONES POR PREGUNTA ═══════════════════════════════════════════════
\section*{Observaciones por pregunta}

\begin{longtable}{>{\bfseries}p{2.8cm} p{1.6cm} p{7.0cm} p{2.8cm}}
  \toprule
  \textbf{Pregunta} & \textbf{Puntaje} & \textbf{Evaluación} & \textbf{Errores detectados} \\
  \midrule
  \endhead
  \bottomrule
  \endfoot
    \textbf{Pregunta 1: Termorregulación y Eficiencia de la Sudoración} & 2.0/2.0 & Excelente respuesta. El estudiante demuestra una comprensión clara de la termorregulación por sudoración, diferenciando la eficiencia entre sudor que gotea y sudor que se evapora. La explicación del efecto de la humedad relativa alta es precisa y bien justificada, incluyendo las consecuencias clínicas. El cálculo es correcto y las unidades son adecuadas. Se detectó un error de transcripción en los datos (escribió 350g pero usó 0.15kg, que es el valor correcto del problema), pero no afectó el resultado final. & \textit{Error de transcripción en la masa (escribió 350g en lugar de 150g, aunque usó el valor correcto de 0.15kg en el cálculo).} \\
    \midrule
    \textbf{Pregunta 2: Mecánica Respiratoria y Síndrome de Dificultad Respiratoria} & 1.4/2.0 & La respuesta identifica correctamente la Ley de Laplace y la relación inversa entre presión y radio. Sin embargo, la explicación de las consecuencias de la tensión superficial constante (flujo de aire y colapso de alvéolos pequeños) es incompleta. La función del surfactante se describe de forma muy concisa, sin detallar cómo iguala las presiones para prevenir el colapso. El cálculo de la diferencia de presión es correcto. & \textit{Explicación incompleta de las consecuencias de la tensión superficial constante en los alvéolos (falta mencionar el flujo de aire y colapso).; Explicación concisa de la función del surfactante, sin detallar el mecanismo de igualación de presiones.} \\
    \midrule
    \textbf{Pregunta 3: Optimización en Hemodiálisis y Primera Ley de Fick} & 1.5/2.0 & La aplicación de la Primera Ley de Fick para calcular el cambio en la eficiencia de difusión es correcta y bien justificada. Se identifica correctamente la fragilidad estructural como riesgo al reducir el espesor de la membrana, pero se omiten las graves consecuencias clínicas (ruptura, mezcla de fluidos, hemólisis, infección). El cálculo de la tasa de transporte es correcto, aunque no se realizó la conversión final a mg/s como se podría esperar. & \textit{Omisión de las consecuencias clínicas graves de la ruptura de la membrana en hemodiálisis (hemólisis, infección).; No se convirtió la tasa de transporte a mg/s al final del cálculo.} \\
    \midrule
    \textbf{Pregunta 4: Errores de Infusión Intravenosa y Ósmosis} & 1.5/2.0 & La identificación de hemólisis y crenación es correcta. La explicación del mecanismo de hemólisis es adecuada. Sin embargo, la descripción del mecanismo de crenación es conceptualmente errónea al indicar que el agua sale 'al interior de los eritrocitos' cuando debería salir *del* interior de los eritrocitos hacia el medio hipertónico. El cálculo de la presión osmótica es correcto. & \textit{Error conceptual en la descripción del mecanismo de crenación: se indica que el agua sale 'al interior de los eritrocitos' en lugar de 'del interior de los eritrocitos'.} \\
    \midrule
    \textbf{Pregunta 5: Capilaridad y Toma de Muestras Sanguíneas} & 0.0/2.0 & La pregunta quedó en blanco, indicando una falta de conocimiento o comprensión sobre el tema de capilaridad. & \textit{Pregunta sin responder.} \\
    \midrule
\end{longtable}

% ══ ERRORES DETECTADOS ═══════════════════════════════════════════════════════
\section*{Análisis de errores}

\subsection*{Errores conceptuales}
\begin{itemize}[leftmargin=1.5em, itemsep=2pt]
  \item Explicación incompleta de las consecuencias de la tensión superficial constante en los alvéolos (Pregunta 2a).
  \item Explicación incompleta de la función del surfactante (Pregunta 2b).
  \item Omisión de las consecuencias clínicas graves de la ruptura de la membrana en hemodiálisis (Pregunta 3b).
  \item Error en la descripción del mecanismo de crenación, invirtiendo la dirección del flujo de agua (Pregunta 4b).
  \item Ausencia total de respuesta en la Pregunta 5, lo que sugiere una falta de comprensión de los principios de capilaridad.
\end{itemize}

\subsection*{Errores procedimentales}
\begin{itemize}[leftmargin=1.5em, itemsep=2pt]
  \item Error de transcripción en los datos de la masa en Pregunta 1c (escribió 350g en lugar de 150g, aunque usó el valor correcto en el cálculo).
  \item No se realizó la conversión final de unidades a mg/s en el cálculo de la tasa de transporte (Pregunta 3c).
\end{itemize}

% ══ RECOMENDACIONES AL ESTUDIANTE ════════════════════════════════════════════
\section*{Retroalimentación para el estudiante}
\begin{tcolorbox}[cajaRecomendacion, title={Dirigido a: Debora Echezuria}]
Estimado estudiante, su examen muestra un buen dominio de los cálculos y la aplicación de fórmulas en varios temas. Sin embargo, es crucial que profundice en la *explicación cualitativa* de los fenómenos físicos y sus *implicaciones biológicas y clínicas*. En varias preguntas, identificó correctamente el fenómeno o la ley, pero la justificación o la descripción del mecanismo subyacente fue incompleta o, en un caso, conceptualmente errónea (como en la crenación). La física en bioanálisis no es solo aplicar fórmulas, sino entender *por qué* y *cómo* ocurren los procesos en el cuerpo. Le recomiendo revisar a fondo los temas de mecánica respiratoria, las consecuencias de la difusión en hemodiálisis y, especialmente, el tema de capilaridad, ya que la Pregunta 5 quedó sin responder. Practique describir los procesos con sus propias palabras, asegurándose de que la lógica física y las consecuencias biológicas sean claras y completas.
\end{tcolorbox}

% ══ NOTA PARA EL PROFESOR ════════════════════════════════════════════════════
\section*{Nota para el profesor}
\begin{tcolorbox}[
  enhanced, breakable,
  colback=yellow!8, colframe=yellow!60!black,
  fonttitle=\bfseries, coltitle=black,
  title={Observaciones para revisión manual},
  top=4pt, bottom=4pt, left=8pt, right=8pt,
]
El estudiante demuestra una base sólida en la aplicación matemática de los principios físicos. Sin embargo, la rúbrica enfatiza la 'comprensión profunda y el análisis físico aplicado a fenómenos biológicos y clínicos', y en este aspecto, el estudiante presenta deficiencias. Las explicaciones cualitativas son a menudo concisas o incompletas, y en un caso (Q4b), conceptualmente incorrectas en la descripción del flujo osmótico. La omisión completa de la Pregunta 5 es un punto débil significativo. Se sugiere reforzar la importancia de la argumentación y el análisis físico-biológico más allá de los cálculos.
\end{tcolorbox}

% ══ PIE DE INFORME ════════════════════════════════════════════════════════════
\vspace{12pt}
\begin{center}
  \footnotesize\color{gray}
  Informe generado automáticamente por el sistema de evaluación con IA (gemini-2.5-flash) y soporte LaTex. \\
  Este documento es una evaluación \textbf{preliminar} para ser revisado por el profesor antes de ser entregado al 
  estudiante. \\
  Generado el 2026-08-08 \textbullet\ BIOANALISIS Sistema Académico de Evaluación.
  \end{center}

\end{document}
